\section{Introduction}\label{sec:intro}

Applications built on large language models are increasingly \emph{workflows}: graphs of model calls, tools and glue code in which one call's output decides what happens next. A router classifies a support ticket and dispatches it to a billing or a technical chain; a grader decides whether retrieved documents are relevant and, if not, triggers a web search; an evaluator rejects a draft and sends it back for another round. The LangGraph documentation teaches these patterns as its core workflow idioms, and Prompt Flow encodes them as conditionally activated nodes.

Every call such a workflow makes leaves a usage record. Providers bill for input and output tokens per call; gateways log the served model and token counts; observability platforms export per-call spans. Operators who handle sensitive content are told to redact it: LangSmith can hide inputs and outputs while ``the trace produced will have its metadata present''~\cite{langsmith2026mask}; Helicone's omit-logs mode drops request bodies but ``you can still see the key metrics for each request''---cost, latency and token usage~\cite{helicone2026omit}; the OpenTelemetry GenAI conventions make message content opt-in but recommend recording the model and the token counts of every call~\cite{otel2026genai}. The redacted trace is then shared with an observability vendor, a gateway operator, a finance team or a security information system.

This paper starts from a simple observation: \emph{when control flow depends on content, the content-free usage trace is a function of the content}. A support ticket routed to human review makes one call, a ticket routed to drafting makes two; a question the private corpus cannot answer triggers a rewrite and another grading round; a ``medical'' request is answered by a longer chain than a ``billing'' one. None of this is hidden by encryption or redaction, because the leak is in the \emph{shape} of the computation, not in its bytes. Recent measurements confirm the danger in adjacent settings: logged router decisions reveal sensitive request topics on real traffic~\cite{chen2026routing}, encrypted traffic fingerprints agent workflows~\cite{zhang2025agentprint,pouryousef2026benchmark}, and an LLM classifier's output length reveals its class~\cite{zhang2024timewilltell}. Agent security frameworks based on information-flow control explicitly leave this channel open: FIDES notes that ``an adversary that is able to see the sequence of tool calls may still be able to infer'' secrets~\cite{costa2025fides}, and CaMeL lists side channels as an open problem~\cite{debenedetti2025camel}.

\paragraph{How common is the problem?} We coded every distinct application workflow in the official LangGraph, LangChain documentation and Prompt Flow examples at pinned revisions (\cref{sec:census}). Of \censusApps{} application workflows, \censusDependent{} (\censusPct\%) branch, loop or fan out on content, and in \censusUserDriven{} of those the decision depends on what the user wrote. In \censusUnbounded{} of \censusLoops{} looping workflows the loop has no iteration bound at all. The channel is not an exotic corner case; it is the default behavior of the idioms these frameworks teach.

\paragraph{Why existing mitigations do not suffice.} Deployed mitigations pad streamed packets with random bytes~\cite{cloudflare2024padding,openai2026obfuscation}; they do not touch usage fields, which are logged unpadded, and random padding leaves measurable residual leakage even at the network level~\cite{mcdonald2025whisper,zhang2025netecho,li2026merged}. Padding a single call's output to its maximum length~\cite{zhang2024timewilltell} hides that call's length but not \emph{which} calls a workflow made, how many, or with which model. Constant-time programming has long solved the analogous problem for CPU instructions by transforming secret branches so that both execute the same operations~\cite{agat2000transforming,molnar2005pc,cauligi2019fact,ngo2017constant}, but LLM calls are not instructions: their cost is set by a tokenizer and by stochastic generation, the observer sees per-call token counts, and a ``dummy'' operation is a billed model call.

\paragraph{This paper.} We present \sys, a language and compiler that makes the usage trace of an LLM workflow independent of its secrets, together with machine-checked guarantees and an evaluation on real models.
\begin{itemize}
\item \textbf{Measurement.} A census of \censusRows{} example workflows and faithful ports of nine representative ones---routing, prompt chaining, evaluator--optimizer loops, orchestrator--worker fan-out, a Prompt Flow switch, an email-support agent, corrective and adaptive retrieval-augmented generation, and a calculator agent---run with Llama~3.2 models over public datasets. From content-free usage records alone, a classifier recovers the workflow's internal secret decision in all nine workflows (\cref{sec:eval-leak}).
\item \textbf{Detection.} \sys labels data with a two-point confidentiality lattice, infers which calls run under secret control or send secret-derived requests, and reports each construct through which a secret can change the usage trace (\cref{sec:language}).
\item \textbf{Repair.} A shaping transformation (\cref{sec:shaping}) pads secret-dependent requests to a tokenizer-exact static bound, forces their outputs to a fixed length, aligns the call sequences of secret branches into a minimum-cost common supersequence filled with dummy calls of the matching model, and pads secret-guarded loops to their bound. Only calls whose usage can depend on a secret are padded.
\item \textbf{Verification.} The compiler is untrusted. A small validator, proved sound in Lean~4 without additional axioms, accepts a shaped program only if it satisfies a padding discipline and only adds padding to the original program (\cref{sec:formal}). We prove that accepted programs emit the same usage trace for all secret inputs, for every randomness tape and therefore in distribution, and compute the original program's results whenever padding is transparent. Both assumptions are stated as explicit hypotheses and tested empirically.
\item \textbf{Evaluation.} Shaped programs collapse each workflow to a single observable trace, so attacks fall to the majority rate, at a median token overhead of \resOverheadMedian$\times$ and with task accuracy within \resFuncMaxDrop{} points of the undefended workflow (\cref{sec:eval}). Ablations corresponding to deployed or proposed defenses---random padding at equal budget, per-call padding without alignment~\cite{zhang2024timewilltell}, and input-only padding on APIs that cannot force output length---all leak. Fault injection shows that the validator catches every leaking compiler bug we injected.
\end{itemize}

\paragraph{What is and is not claimed.} The guarantee concerns an observer of per-call usage records: model identity, input and output token counts, and tool names. It is exact for that observer under the stated conformance assumptions, which self-hosted serving systems satisfy through exact tokenizers and forced-length decoding (\texttt{ignore\_eos}, \texttt{min\_tokens})~\cite{vllm2026sampling,llamacpp2026server}. Wall-clock latency is not covered by the theorem; we measure it separately (\cref{sec:eval-latency}). Hosted APIs that cannot force output length support only a weaker input-and-shape padding whose residual leakage we quantify. Padding costs tokens, and we report that cost honestly; the transformation pays it only where secrets can flow.

All code, proofs, workflow ports, datasets, raw run records and analysis scripts are released with the paper (\cref{sec:artifact}).
